\textbf{Polarizace} je posun vázaných nábojů, orientace dipólů nebo hromadění nábojů na nehomogenitách působením elektrického pole. Čtyři mechanismy od nejrychlejšího:
\begin{enumerate}
\item \textbf{Elektronová:} posun elektronového obalu vůči jádru, $10^{-16}$--$10^{-14}$ s.
\item \textbf{Iontová pružná:} vzájemný posun kladných a záporných iontů v mřížce, $10^{-13}$--$10^{-12}$ s.
\item \textbf{Dipólová orientační:} natáčení permanentních dipólů do směru pole; u nízkomolekulárních látek $10^{-12}$--$10^{-8}$ s.
\item \textbf{Mezivrstvová:} nahromadění pohyblivých nábojů na rozhraních nehomogenit; z uvedených nejpomalejší.
\end{enumerate}
